\section{Metodología de recolección y limpieza}\label{sec:metodologia}

\subsection{Organización del proyecto}
Los datos se organizan en tres niveles. Así, los archivos originales nunca se modifican y cada
transformación puede reproducirse:
\begin{center}
\begin{tabular}{ll}
\toprule
\textbf{Carpeta} & \textbf{Contenido}\\
\midrule
\archivo{data/raw/}       & Descargas originales, sin modificar \\
\archivo{data/interim/}   & Datos limpios por fuente (CSV y Parquet)\\
\archivo{data/processed/} & Base integrada (alcaldía × mes, CDMX × semana)\\
\archivo{docs/}           & Diccionarios de datos y notas metodológicas\\
\archivo{src/download/}   & Un script de descarga por fuente\\
\archivo{src/clean/}      & Un script de limpieza por fuente\\
\archivo{notebooks/}      & Análisis exploratorio\\
\archivo{informe/}        & Este documento\\
\bottomrule
\end{tabular}
\end{center}

\subsection{Criterios de limpieza}
Estos criterios aplican a todas las fuentes:
\begin{enumerate}
  \item \textbf{Codificación.} Todos los archivos se convierten a UTF-8. Por ejemplo, el CSV del Metro
        presenta texto mal codificado (\emph{mojibake}): \texttt{Isabel la Católica} en lugar de
        \enquote{Isabel la Católica}.
  \item \textbf{Fechas.} Se convierten al formato ISO~8601. En las carpetas de la FGJ hay que decidir
        explícitamente si se usa la fecha del hecho o la fecha de inicio de la carpeta.
  \item \textbf{Alcaldías.} Se crea un catálogo único con la clave INEGI (\archivo{09002}--\archivo{09017}).
        Se normalizan acentos, mayúsculas y nombres anteriores (\enquote{Delegación}).
  \item \textbf{Estaciones y líneas.} Se crea un catálogo único y se documentan los cierres prolongados:
        la Línea~12 desde el 3 de mayo de 2021 y la Línea~1 por su modernización desde julio de 2022.
  \item \textbf{Faltantes y ceros.} Se distingue un día \enquote{sin servicio} (cierre) de un \enquote{dato faltante}.
  \item \textbf{Duplicados y registros fuera del dominio.} Se eliminan los duplicados y las carpetas
        con coordenadas fuera del polígono de la CDMX.
  \item \textbf{Catálogo de delitos.} Los delitos se agrupan en categorías estables (alto impacto,
        patrimonial, violencia familiar, etc.) que no cambien con las reclasificaciones de la FGJ.
  \item \textbf{Agregación.} Toda serie se entrega también a nivel alcaldía × mes y CDMX × semana.
  \item \textbf{Trazabilidad.} De cada archivo descargado se registran la URL, la fecha de descarga
        y el hash SHA-256 en \archivo{data/raw/MANIFEST.csv}.
\end{enumerate}

\subsection{Integración}
Las llaves comunes son la \textbf{fecha} (día, semana o mes) y la \textbf{clave de alcaldía}.
Cada fuente limpia se agrega a esas llaves y después se unen todas. Para comparar series con escalas
distintas se calculan dos indicadores:
\begin{itemize}
  \item \textbf{Tasas por cada 100 mil habitantes}, para los delitos.
  \item \textbf{Índices base 100}, donde 100 es el promedio de 2019. Así se mide la caída y la
        recuperación en términos relativos.
\end{itemize}
